R01 & Prise de contrôle d'un compte à privilèges menant à une modification non autorisée de l'infrastructure (AD, serveurs) & AD01, BASTION01 et VPN, groupe Server-Admins & 4 comptes privilégiés sur 10 sans MFA dont admin-old dormant et admin-presta, RDP direct FW-17, AL-1048 non traitée & 4 & 3 & \textbf{12} & Réduire (A08, A02) et mesures conservatoires immédiates. Décision DG sur proposition du RSSI. \\
R02 & Indisponibilité de MEDI-PORTAL au-delà du RTO de 8 h faute de restauration fiable (rançongiciel ou panne) & MEDI-PORTAL, base de données MEDI, BACKUP01 & 2 tests documentés sur 9 services, RPO inconnu, FW-31 Any des postes vers les sauvegardes, svc-backup non restreint & 4 & 3 & \textbf{12} & Réduire (A05, A08) et restriction de FW-31 en régie. Décision DG. \\
R03 & Altération ou divulgation de données patients via l'application métier non corrigée & APP01, DB01 & Correctif critique manquant depuis 74 jours, aucun SLA, flux WEB01 vers APP01, RDP FW-17 & 4 & 2 & \textbf{8} & Réduire (A09). Exception signée et mesures compensatoires si le correctif est reporté. \\
R04 & Divulgation de données de santé depuis les partages de fichiers après compromission d'un poste & SRV-FICHIERS, AD01 & SMB ouvert à tout le VLAN 40 (FW-19), signature non obligatoire, retard de correctifs de 43 jours & 3 & 3 & \textbf{9} & Réduire partiellement (A09, A03). Segmentation complète A10 différée, risque résiduel accepté par la DG pour 12 mois. \\
R05 & Interception ou dégradation des sessions du portail public & WEB01, MEDI-PORTAL & TLS 1.0 signalé, port 80 ouvert, redirection non vérifiée & 2 & 2 & \textbf{4} & Réduire à faible coût dans A09, sinon accepter temporairement. \\
R06 & Indisponibilité ou altération de service causée par un changement mal maîtrisé & MEDI-PORTAL, infrastructure & Conception et réalisation non standardisées (P13, P14), règle provisoire non retirée (FW-17) & 3 & 3 & \textbf{9} & Accepter temporairement (A04 différée en année 2) avec revue de changement renforcée en régie. Risque résiduel porté par la DO. \\
